\chapter{Complex Functions} % 复变函数
\section{Complex Functions and Analytic Functions} % 复函数与解析函数 
\begin{leftbarTitle}{Complex Functions}\end{leftbarTitle}
\begin{definition}{Complex Functions}
    A \textbf{complex function} is a function 
    \[
    f: 
    \begin{matrix}
        E \to \mathbb{C}, \\
        z \mapsto \omega,
    \end{matrix}
    \] 
    where \(E \subseteq \mathbb{C}\) is the domain of \(f\) and \(\omega = f(z)\) is the value of \(f\) at \(z\).
    
    % 单值函数与多值函数
    A complex function is called \textbf{single-valued} if 
    for every \(z \in E\), there is a unique value \(\omega\) such that \(f(z) = \omega\).
    Otherwise, it is called \textbf{multi-valued}.
\end{definition}

\begin{caution}
    % 需要注意的是, 所谓的多值函数并不是一个真正意义上的函数, 因为它不满足函数的定义, 但我们仍然习惯性地称其为函数, 以便于讨论和分析
    % 之后我们会通过解析延拓等方法将多值函数转化为单值函数, 从而使其成为一个真正意义上的函数
    It should be noted that a so-called multi-valued function \emph{is not a function in the true sense}, 
    because it does not satisfy the definition of a function. 
    However, we still habitually call it a function for the sake of discussion and analysis.
    Later, we will use methods such as analytic continuation to transform multi-valued functions into 
    single-valued functions, thus making them true functions in the mathematical sense.
\end{caution}





\begin{leftbarTitle}{Limit, Continuity, and Differentiability}\end{leftbarTitle}
\begin{definition}{Limit of Complex Functions}
    Let \(f: E \to \mathbb{C}\) be a complex function and \(z_0\) be a limit point of \(E\). 
    We say that the limit of \(f(z)\) as \(z\) approaches \(z_0\) is \(\omega_0\), denoted by
    \[
    \lim_{z \to z_0} f(z) = \omega_0,
    \]
    if for every \(\varepsilon > 0\), there exists a \(\delta > 0\) such that 
    whenever \(0 < |z - z_0| < \delta\) and \(z \in E\), we have \(|f(z) - \omega_0| < \varepsilon\).
\end{definition}

\begin{definition}{Continuity of Complex Functions}
    A complex function \(f: E \to \mathbb{C}\) is said to be \textbf{continuous} at a point \(z_0 \in E\) if 
    \[
    \lim_{z \to z_0} f(z) = f(z_0).
    \]
    If \(f\) is continuous at every point in its domain \(E\), then we say that \(f\) is continuous on \(E\).
\end{definition}

\begin{definition}{Uniform Continuity of Complex Functions}
    A complex function \(f: E \to \mathbb{C}\) is said to be \textbf{uniformly continuous} on \(E\) if 
    for every \(\varepsilon > 0\), there exists a \(\delta > 0\) such that 
    whenever \(z_1, z_2 \in E\) and \(|z_1 - z_2| < \delta\), we have \(|f(z_1) - f(z_2)| < \varepsilon\).
\end{definition}

\begin{definition}{Derivative and Differential of Complex Functions}
    Let \(f: E \to \mathbb{C}\) be a complex function and \(z_0 \in E\). 
    We say that \(f\) is \textbf{derivable} at \(z_0\) if the limit
    \[
    f'(z_0) = \lim_{z \to z_0} \frac{f(z) - f(z_0)}{z - z_0}
    \]
    exists. The value \(f'(z_0)\) is called the \textbf{derivative} of \(f\) at \(z_0\).
    
    The \textbf{differential} of \(f\) at \(z_0\) is defined as
    \[
    \mathrm{d}f(z_0) = f'(z_0) \mathrm{d}z,
    \]
    where \(\mathrm{d}z\) is an infinitesimal change in \(z\).
\end{definition}

\begin{note}
    % 需要注意的是, 复函数的可微与可导是等价的, 这与多元实函数不同, 因为复函数的导数必须满足 Cauchy-Riemann 方程, 这使得复函数的导数具有更强的性质, 比如解析性和共形性等
    It should be noted that for complex functions, differentiability and derivability are equivalent. 
    This is different from multivariable real functions, where derivability does not necessarily imply differentiability. 
    For complex functions, the derivative must satisfy the Cauchy-Riemann equations, 
    which endow the derivative with stronger properties such as analyticity and conformality.
\end{note}

\begin{leftbarTitle}{Analytic Functions}\end{leftbarTitle}
\begin{definition}{Analytic Functions}
    A complex function \(f: E \to \mathbb{C}\) is said to be \textbf{analytic} at a point \(z_0 \in E\) if 
    there exists an open neighborhood \(U\) of \(z_0\) such that \(f\) is derivable at every point in \(U\).
    If \(f\) is analytic at every point in its domain \(E\), then we say that \(f\) is analytic on \(E\).

    An analytic function is also called a \textbf{holomorphic function}.
\end{definition}

\begin{theorem}{Cauchy-Riemann Equations}
    Let \(f: E \to \mathbb{C}\) be a complex function, where \(E \subseteq \mathbb{C}\) is an open set. 
    Write \(f(z) = u(x, y) + iv(x, y)\), where \(z = x + iy\), 
    and \(u\) and \(v\) are real-valued functions of two real variables \(x\) and \(y\). 
    Then \(f\) is analytic at a point \(z_0 = x_0 + iy_0 \in E\) 
    if and only if the following \textbf{Cauchy-Riemann equations} hold at \((x_0, y_0)\):
    \[
    \begin{cases}
        \frac{\partial u}{\partial x}(x_0, y_0) = \frac{\partial v}{\partial y}(x_0, y_0), \\
        \frac{\partial u}{\partial y}(x_0, y_0) = -\frac{\partial v}{\partial x}(x_0, y_0).
    \end{cases}
    \]
\end{theorem}

\section{Elementary Functions} % 初等函数
\begin{definition}{Exponential Functions}
    Let \(z=x+\mathrm{i}y\) with \(x, y \in \mathbb{R}\). The exponential function \(f(z) = e^z\) is defined as
    \[
    e^z = e^x (\cos y + \mathrm{i} \sin y),
    \]
    which satisfies:
    \begin{enumerate}
        \item \(\forall x\in \mathbb{R}, f(x)=e^x\); 
        \item \(f(z)\) is analytic on \(\mathbb{C}\);
        \item \(\forall z_1, z_2 \in \mathbb{C}, f(z_1 + z_2) = f(z_1) f(z_2)\).
    \end{enumerate}
\end{definition}

\begin{property}
    \begin{enumerate}
        \item \(f(z)\) is periodic with period \(2\pi \mathrm{i}\), i.e., \(\forall z \in \mathbb{C}, f(z + 2\pi \mathrm{i}) = f(z)\);
        \item \(f(z)\) is never zero, i.e., \(\forall z \in \mathbb{C}, f(z) \neq 0\);
        \item \(f(z)\) is unbounded, i.e., \(\forall M > 0, \exists z \in \mathbb{C}\) such that \(|f(z)| > M\).
    \end{enumerate}
\end{property}


\vspace{0.7cm}
\begin{definition}{Power Functions and Root Functions}
    Let \(\alpha \in \mathbb{C}\) be a fixed complex number. The power function \(f(z) = z^\alpha\) is defined as
    \[
    z^\alpha = e^{\alpha \mathrm{Ln}\, z},
    \]
    where \(\mathrm{Ln}\, z\) is the logarithmic function defined above.
    When \(z=0\), define \(z^{\alpha} =0\).
    
    The root function \(f(z) = \sqrt[n]{z}\) is defined as
    \[
    \sqrt[n]{z} = z^{\frac{1}{n}} = e^{\frac{1}{n} \mathrm{Ln}\, z},
    \]
    where \(n\) is a positive integer.
\end{definition}


\vspace{0.7cm}
\begin{definition}{Trigonometric Functions}
    The trigonometric functions \(\sin z\) and \(\cos z\) are defined as
    \begin{gather*}
        \sin z = \frac{e^{\mathrm{i} z} - e^{-\mathrm{i} z}}{2\mathrm{i}}, \\
        \cos z = \frac{e^{\mathrm{i} z} + e^{-\mathrm{i} z}}{2}.  
    \end{gather*}
\end{definition}

\begin{property}
    \begin{enumerate}
        \item \(\cos z\) is even and \(\sin z\) is odd, i.e., 
            \(\forall z \in \mathbb{C}, \cos(-z) = \cos z\) and \(\sin(-z) = -\sin z\);
        \item \(\sin z\) and \(\cos z\) are periodic with period \(2\pi\), i.e., 
            \(\forall z \in \mathbb{C}, \sin(z + 2\pi) = \sin z\) and \(\cos(z + 2\pi) = \cos z\);
        \item \(\sin z\) and \(\cos z\) are unbounded, i.e., 
            \(\forall M > 0, \exists z \in \mathbb{C}\) such that \(|\sin z| > M\) and \(|\cos z| > M\).
        \item \(\sin z\) and \(\cos z\) are analytic on \(\mathbb{C}\) and 
            \(\left( \sin z \right)' = \cos z\), \(\left( \cos z \right)' = -\sin z\).
    \end{enumerate}  
\end{property}

\section{Elementary Multi-Valued Functions} 
\begin{definition}{Branch Points, Branch Cuts, and Branches} % 支点、割线与分支
    Let \(f\) be a multi-valued complex function on a domain \(E \subseteq \widehat{\mathbb{C}}\), 
    where \(\widehat{\mathbb{C}}\) denotes the Riemann sphere.
    
    A point \(z_0 \in E\) is called a \textbf{branch point} of \(f\) if the analytic continuation of a local branch of \(f\)
    along a closed loop encircling \(z_0\) does not return to its initial value.
    
    A \textbf{branch cut} of \(f\) is a curve (or a set of curves) removed from \(E\) so that, on the remaining domain,
    one can choose a single-valued branch of \(f\) (typically on a simply connected domain).

    A \textbf{branch} of \(f\) is a single-valued function defined on a domain obtained 
    by cutting the complex plane along the branch cuts,
    such that it coincides with \(f\) on that domain.

    % 单值连续分支与单值解析分支
    A branch can be further classified into a \textbf{single-valued continuous branch} if it is continuous,
    and an \textbf{analytic branch} if it is analytic.
\end{definition}


\begin{definition}{Argument Function}
    Let \(z\in \mathbb{C} \setminus \{0\}\) be a nonzero complex number. 
    The \textbf{argument function} \(\arg(z)\) is defined as:
    \[
    \omega = \mathrm{Arg} (z) = \arg(z) + 2k\pi, \quad k \in \mathbb{Z},
    \]
    where \(\arg(z)\) is the principal value of the argument, which takes values in \((-\pi, \pi]\).
\end{definition}



For the argument function, the branch points are \(0\) and \(\infty\): 
a circuit around either point changes \(\mathrm{Arg}\, z\) by an integer multiple of \(2\pi\), 
so monodromy is nontrivial. 
A standard branch cut is the negative real axis \(( -\infty, 0 ]\), 
and on \(\mathbb{C} \setminus ( -\infty, 0 ]\) one obtains the principal branch
\[
\arg z \in (-\pi, \pi).
\]
This branch is single-valued and continuous, but not analytic (indeed, real-valued and non-constant on a planar domain).


\vspace{0.7cm}
\begin{definition}{Logarithmic Functions}
    The logarithmic function \(f(z) = \mathrm{Ln}\, z\) is defined as the inverse function of the exponential function \(e^z\), i.e., 
    \[
    \omega = f(z) = \mathrm{Ln}\, z \iff e^{\omega} = z.
    \]
    Since the exponential function is periodic with period \(2\pi \mathrm{i}\), the logarithmic function is multi-valued, 
    and can be expressed as
    \[
    f(z) = \mathrm{Ln}\, z = \ln |z| + \mathrm{i} \mathrm{Arg}\, z,
    \]
    where \(\ln |z|\) is the natural logarithm of the modulus of \(z\).
\end{definition}
% 由此可见, 任何非零复数有无穷多个对数, 任意两个相差 \(2k\pi \mathrm{i}\), \(k \in \mathbb{Z}\).
It can be seen that \emph{any nonzero complex number has infinitely many logarithms,
and any two of them differ by \(2k\pi \mathrm{i}\) for some integer \(k\)}.

% 对数函数的支点、割线与分支
The branch points of the logarithmic function are \(0\) and \(\infty\),
and a standard branch cut is the negative real axis \(( -\infty, 0 ]\).
On \(\mathbb{C} \setminus ( -\infty, 0 ]\), one can define the principal branch of the logarithmic function as
\[
\mathrm{ln}\, z = \ln |z| + \mathrm{i} \arg z,
\]
which is single-valued and analytic, so it is an analytic branch of the logarithmic function.